\cotefile{12}
\scanpage{2}
% archive foliation: 1 (upper right; also lower left)
\uline{Esquisse d'une théorie des motifs.}

1. Revue des théories cohomologiques habituelles des schémas projectifs lisses.
L'idée de motifs, \inserted{\unsure{qui doit remplacer|que doit remplir} la cohomologie \unsure{entière} (en terme des modules de \unsure{t.f.|déf.} sur $\mathbb{Z}$ !).}
\unsure{Réserves} : a) Isomotifs, b) Motifs $\frac12$ simples c) Base ponctuelle.
\margin{*}

2. Définition de la catégorie $\mathcal{M}^+(k)$ des motifs effectifs sur $k$.
C'est une catégorie abélienne \inserted{``$\frac12$ simple''} avec $\otimes$.
Foncteurs $T^{\unsure{*}}_{\mathrm{DR}}$, $T_{\mathrm{Hdg}}$ (si $k$ de car nulle), $T_B$ (si $k=\mathbb{C}$), $T^\ell$ (si $\ell\ne\deleted{\unsure{car}}\inserted{\operatorname{car}}k$) \deleted{\illegible{3}} ces foncteurs sont multiplicatifs. Tous ces foncteurs sont fidèles.

3. Graduation dans la catégorie $\mathcal{M}^+(k)$. Les foncteurs $h^i(X)$.\unsure{$\mathcal{M}^{+i}(k)$}.
\inserted{Le produit $\otimes$ compatible avec graduation.}

4. L'équivalence $\mathcal{M}^{+1}(k)^\circ\simeq\operatorname{Isab}(k)$
($X\mapsto\underline{\operatorname{Pic}}^0_X$, $A\mapsto h^1(\widehat A)$).

\deleted{4. Le motif $\mathbb{Q}(-1)$, et la catégorie $\mathcal{M}(k)$ des motifs \unsure{et}}
\margin{*}
Les motifs de poids $\geq1$ comme généralisation des VA. Le \unsure{caractère} ``discret'' est conservé (s'exprimant \unsure{alg.ment} par le fait que si $k\longrightarrow k'$ est une extension de corps algébriquement clos, alors $\mathcal{M}^{\unsure{e}}(k)\unsure{-|\longrightarrow}\mathcal{M}^{\unsure{e}}(k')$ est pleinement fidèle).

\unsure{Caractère} continu\dots\ Exemples d'énoncé suggéré par la théorie des \deleted{schémas} abéliens : caractérisation des motifs \uuline{lisses} sur $S$ \deleted{\unsure{géométriquement réduit}}\inserted{normal}\deleted{\unsure{et connexe}}.
\margin{*}

5. Le motif $\mathbb{Q}(\unsure{-1|+1})$ (anti-Tate). La catégorie $\mathcal{M}(k)$ des motifs (isomotifs semi-simples sur $k$ !). Maintenant il y a des \uline{Hom} internes, en plus des $\otimes$ et de la graduation.

Si $\mathcal C$ est une catégorie pseudo-abélienne avec $\otimes$, alors les foncteurs compatibles avec $\otimes$ de $\mathcal{M}(k)$ dans $\mathcal C$ correspondent aux foncteurs $T:V(k)^\circ\longrightarrow\mathcal C$, satisfaisant Künneth, et tels que \deleted{\unsure{$\mathbb{P}^1_k$}} $T(\mathbb{Q}(-1))=\operatorname{Im}T(\pi)$
($\pi$ le projecteur dans $h(\mathbb{P}^1_k)$ \deleted{\illegible{1}}, d'im.\ $h^2(\mathbb{P}^1_k)=\mathbb{Q}(-1)$)
\uline{soit inversible}. On montre qu'un tel $T$ commute automatiquement à la formation des \uline{Hom}. La catégorie des motifs $\mathcal{M}(k)$ est donc caractérisée, avec sa structure $\otimes$, \deleted{et l'objet $\mathbb{Q}(\unsure{-1})$} à équivalence près par propriété universelle.
